\documentclass[10pt, letterpaper]{article}

% Packages:
\usepackage[
    ignoreheadfoot, % set margins without considering header and footer
    top=1 cm, % seperation between body and page edge from the top
    bottom=2 cm, % seperation between body and page edge from the bottom
    left=2 cm, % seperation between body and page edge from the left
    right=2 cm, % seperation between body and page edge from the right
    footskip=1.0 cm, % seperation between body and footer
    % showframe % for debugging 
]{geometry} % for adjusting page geometry
\usepackage{titlesec} % for customizing section titles
\usepackage{tabularx} % for making tables with fixed width columns
\usepackage{array} % tabularx requires this
\usepackage[dvipsnames]{xcolor} % for coloring text
\definecolor{primaryColor}{RGB}{0, 0, 0} % define primary color
\usepackage{enumitem} % for customizing lists
\usepackage{fontawesome5} % for using icons
\usepackage{amsmath} % for math
\usepackage[
    pdftitle={Jiacheng's CV},
    pdfauthor={Jiacheng},
    pdfcreator={LaTeX with RenderCV},
    colorlinks=true,
    urlcolor=primaryColor
]{hyperref} % for links, metadata and bookmarks
\usepackage[pscoord]{eso-pic} % for floating text on the page
\usepackage{calc} % for calculating lengths
\usepackage{bookmark} % for bookmarks
\usepackage{lastpage} % for getting the total number of pages
\usepackage{changepage} % for one column entries (adjustwidth environment)
\usepackage{paracol} % for two and three column entries
\usepackage{ifthen} % for conditional statements
\usepackage{needspace} % for avoiding page brake right after the section title
\usepackage{iftex} % check if engine is pdflatex, xetex or luatex
\usepackage{graphicx}
\usepackage{ragged2e}
\usepackage{tabularx} % 导言区引入

% Ensure that generate pdf is machine readable/ATS parsable:
\ifPDFTeX
    \input{glyphtounicode}
    \pdfgentounicode=1
    \usepackage[T1]{fontenc}
    \usepackage[utf8]{inputenc}
    \usepackage{lmodern}
\fi

\usepackage{charter}
% Some settings:
\raggedright
\AtBeginEnvironment{adjustwidth}{\partopsep0pt} % remove space before adjustwidth environment
\pagestyle{empty} % no header or footer
\setcounter{secnumdepth}{0} % no section numbering
\setlength{\parindent}{0pt} % no indentation
\setlength{\topskip}{0pt} % no top skip
\setlength{\columnsep}{0.15cm} % set column seperation
\pagenumbering{gobble} % no page numbering

\titleformat{\section}{\needspace{4\baselineskip}\bfseries\large}{}{0pt}{}[\vspace{1pt}\titlerule]

\titlespacing{\section}{
    % left space:
    -1pt
}{
    % top space:
    0.3 cm
}{
    % bottom space:
    0.2 cm
} % section title spacing

\renewcommand\labelitemi{$\vcenter{\hbox{\small$\bullet$}}$} % custom bullet points
\newenvironment{highlights}{
    \begin{itemize}[
        topsep=0.10 cm,
        parsep=0.10 cm,
        partopsep=0pt,
        itemsep=0pt,
        leftmargin=0 cm + 10pt
    ]
}{
    \end{itemize}
} % new environment for highlights


\newenvironment{highlightsforbulletentries}{
    \begin{itemize}[
        topsep=0.10 cm,
        parsep=0.10 cm,
        partopsep=0pt,
        itemsep=0pt,
        leftmargin=10pt
    ]
}{
    \end{itemize}
} % new environment for highlights for bullet entries

\newenvironment{onecolentry}{
    \begin{adjustwidth}{
        0 cm + 0.00001 cm
    }{
        0 cm + 0.00001 cm
    }
}{
    \end{adjustwidth}
} % new environment for one column entries

\newenvironment{twocolentry}[2][]{
    \onecolentry
    \def\secondColumn{#2}
    \setcolumnwidth{\fill, 4.5 cm}
    \begin{paracol}{2}
}{
    \switchcolumn \raggedleft \secondColumn
    \end{paracol}
    \endonecolentry
} % new environment for two column entries

\newenvironment{threecolentry}[3][]{
    \onecolentry
    \def\thirdColumn{#3}
    \setcolumnwidth{, \fill, 4.5 cm}
    \begin{paracol}{3}
    {\raggedright #2} \switchcolumn
}{
    \switchcolumn \raggedleft \thirdColumn
    \end{paracol}
    \endonecolentry
} % new environment for three column entries

\newenvironment{header}{
    \setlength{\topsep}{0pt}\par\kern\topsep\centering\linespread{1.5}
}{
    \par\kern\topsep
} % new environment for the header

\newcommand{\placelastupdatedtext}{% \placetextbox{<horizontal pos>}{<vertical pos>}{<stuff>}
  \AddToShipoutPictureFG*{% Add <stuff> to current page foreground
    \put(
        \LenToUnit{\paperwidth-2 cm-0 cm+0.05cm},
        \LenToUnit{\paperheight-1.0 cm}
    ){\vtop{{\null}\makebox[0pt][c]{
        \small\color{gray}\textit{Last updated in September 2024}\hspace{\widthof{Last updated in September 2024}}
    }}}%
  }%
}%

% save the original href command in a new command:
\let\hrefWithoutArrow\href

% new command for external links:

\begin{document}
    \newcommand{\AND}{\unskip
        \cleaders\copy\ANDbox\hskip\wd\ANDbox
        \ignorespaces
    }
    \newsavebox\ANDbox
    \sbox\ANDbox{$|$}

    \begin{header}
        \centering
        \fontsize{25 pt}{25 pt}\selectfont Jiacheng Xue
    
        \vspace{0 pt}
    
        \normalsize
        \mbox{Telephone: +86 187-2968-9067 | Email: jiachengxue2001@gmail.com | \href{https://jcxue001.github.io/personal_website/}{Website} | \href{https://github.com/jcXue001}{GitHub}} 


    \end{header}

    \vspace{5 pt - 0.3 cm}

\section{Education}
\begin{tabularx}{\linewidth}{@{}lX@{\extracolsep{\fill}}r@{}}
    \textbf{Xi'an Jiaotong University} && Shaanxi, China \\
    \vspace{0.2 cm}
    Master of Mechanical Engineering, GPA: 86.21/100 && \textit{Sept. 2023 – present} \\

    \textbf{Xi'an Jiaotong University} && Shaanxi, China \\
    Bachelor of Mechanical Engineering, GPA: 85.61/100 (Rank:18/190) && \textit{Sept. 2019 – Jul. 2023} \\
\end{tabularx}

\section{Research Interests}
\begin{onecolentry}
    \begin{highlightsforbulletentries}
        \item \textbf{Robotics:} Design and control of robotic arms and their reconfigurable, adaptive systems for versatile task environments, including manufacturing, inspection, and human-robot interaction.

        \item \textbf{AI-based Metamaterials Design:} Inverse design of mechanical metamaterials driven by artificial intelligence, with an emphasis on balancing trade-offs between multiple functional and environmental constraints.

        \item \textbf{AI-assisted 3D Printing:} Closed-loop control strategies and online monitoring systems for additive manufacturing processes, enabling defect detection and real-time process correction.
    \end{highlightsforbulletentries}
\end{onecolentry}


\section{Publications}
\begin{onecolentry}
    \newlength{\contentwidth}
    \newcommand{\addDate}[1]{%
        \setbox0=\lastbox 
        \settowidth{\contentwidth}{\usebox0} 
        \unhbox0 
        \hfill\rlap{\textit{\makebox[\contentwidth][r]{#1}}}% 日期斜体
    }
    \begin{highlightsforbulletentries}
        \item \textbf{Xue, J.}, Bao, H., Liu, T., \textit{et.al.} (2025). Inverse design of mechanical metamaterials balancing manufacturability and compactness: A case study on lattice cells. \textit{Extreme Mechanics Letters}, \textit{79}, 102395. \href{https://doi.org/10.1016/j.eml.2025.102395}{https://doi.org/10.1016/j.eml.2025.102395}.


        \item Chi, X.\textsuperscript{†}, \textbf{Xue, J.}\textsuperscript{†}, Jia, L., \textit{et al.} (2025). Machine learning-based online monitoring and closed-loop controlling for 3D printing of continuous fiber-reinforced composites. \textit{Additive Manufacturing Frontiers}, \textit{4}, 200196. \href{https://doi.org/10.1016/j.amf.2025.200196}{https://doi.org/10.1016/j.amf.2025.200196}. \textsuperscript{†}\textit{Co-first authors}.

        \item Wu, L., \textbf{Xue, J.}, Tian, X., \textit{et al.} (2023). 3D-printed metamaterials with versatile functionalities. \textit{Additive Manufacturing Frontiers}, \textit{2}, 100091. \href{https://doi.org/10.1016/j.cjmeam.2023.100091}{https://doi.org/10.1016/j.cjmeam.2023.100091}.

        \item Wu, L., Lu, Y., Li, P., Wang, Y., \textbf{Xue, J.}, \textit{et al.} (2024). Mechanical metamaterials for handwritten digits recognition. \textit{Advanced Science}, \textit{11}, 2308137. \href{https://doi.org/10.1002/advs.202308137}{https://doi.org/10.1002/advs.202308137}.

        \item Song, X., Yan, S., Wang, Y., Zhang, H., \textbf{Xue, J.}, \textit{et al.} (2025). Genetic algorithm-enabled mechanical metamaterials for vibration isolation with different payloads. \textit{Journal of Materiomics}, \textit{11}, 100944. \href{https://doi.org/10.1016/j.jmat.2024.100944}{https://doi.org/10.1016/j.jmat.2024.100944}.
        \end{highlightsforbulletentries}
\end{onecolentry}

\section{Research Experience}

        \begin{twocolentry}{
            \textit{Mar. 2021 – Apr. 2023}
        }
            \textbf{Advanced Mobile Chassis Design}\end{twocolentry}
        \vspace{0.10 cm}
        \begin{onecolentry}
            \begin{highlights}
                \item Engineered a series of innovative mobile chassis for robot systems in demanding environments, incorporating Mecanum wheel and four-steering wheel to enable omnidirectional agility and robust performance.

                \item Pioneered mobility enhancements by tripling robot speed (147\% improvement) and boosting efficiency through precise tuning of motor reduction ratios and center of gravity positioning in Webots simulations, ensuring fluid, stable navigation under competitive pressures.
                
                \item innovated a spherical robot encased in a sealed spherical shell seamlessly , integrating worm drive, gear drive, and belt drive mechanisms to achieve precise, discrete gimbal control, ensuring chassis flexibility and gimbal stability .
                
            \end{highlights}
        \end{onecolentry}
        \vspace{0.2 cm}

        \begin{twocolentry}{
            \textit{Sept. 2022 – Sept 2023}
        }
            \textbf{Online Monitoring and Closed-Loop Control of 3D Printing}\end{twocolentry}
        \vspace{0.10 cm}
        \begin{onecolentry}
            \begin{highlights}
                \item Devised an advanced system integrating Mask-RCNN and Octoprint for real-time fracture detection and parameter monitoring in 3D printing, achieving 94.7\% and 97.7\% recognition accuracies for pure materials and fiber-reinforced composites,
                \item Amplified tensile strength by 17.3 times and elastic modulus by 44.8 times through closed-loop control.
                \item Authored an original research paper synthesizing experimental findings on 3D printing of composites, resulting in an online publication; presented the work at the 13th Asian-Australasian Conference on Composite Materials (ACCM) in Kyoto, Japan.
                
            \end{highlights}
        \end{onecolentry}
        \vspace{0.2 cm}

        \begin{twocolentry}{
            \textit{Oct. 2023 – Feb. 2024}
        }
            \textbf{Pedestrian Trajectory Recognition and Prediction Based on UAV Vision}\end{twocolentry}
        \vspace{0.10 cm}
        \begin{onecolentry}
            \begin{highlights}
                \item Converted the Trajnet dataset into SDD format for YOLOv8 training and employed the ByteTrack tool to effectively track pedestrian trajectories.
                
                \item Analyzed Social-GAN to predict pedestrian trajectories based on their prior locations. By inputting 12 frames of imagery, the process successfully generated 8 frames of predicted trajectories, achieving an ADE (Average Displacement Error) of 0.014.

            \end{highlights}
        \end{onecolentry}
        
        \begin{twocolentry}{
            \textit{Oct. 2023 – present}
        }
            \textbf{Inverse Design of Mechanical Metamaterials Using Machine Learning}\end{twocolentry}
        \vspace{0.10 cm}
        \begin{onecolentry}
            \begin{highlights}
                \item Validated mechanical metamaterial designs through finite element method (FEM) simulations and crafted a bespoke computer vision-based program for non-contact Poisson’s ratio monitoring, significantly enhancing experimental analysis of negative Poisson’s ratio and mechanical response.
                
                \item Analyzed and refined open source machine learning code from research papers to pinpoint the most effective algorithm for inverse design, designed for regression tasks, achieving an exceptional R² accuracy of 99.7\%.
                
                \item Spearheaded the independent design, research, and authorship of a trailblazing paper integrating inverse design with multi-objective optimization for mechanical metamaterials, expertly optimizing lattice structure parameters to balance manufacturability, compactness, and performance.
            \end{highlights}
        \end{onecolentry}

        \begin{twocolentry}{
            \textit{Apr. 2024 – present}
        }
            \textbf{Mobile Robotic Platform for Additive Manufacturing }\end{twocolentry}

        \vspace{0.10 cm}
        \begin{onecolentry}
            \begin{highlights}
                \item Independently constructed a gantry-based mobile robotic platform for additive manufacturing, integrating XYZ moving and rotating axes to enable full degrees of freedom for the 3D printing nozzle, attaining a precision positioning error of 1.41 mm over a 200 mm range.
                
                \item Programmed an STM32 microcontroller board to execute PID algorithms for motor control and streamline communication via CAN, I2C, and UART protocols, discretizing the control system to elevate robotic mobility and printing accuracy.

                \item Secured a Chinese invention patent (No. 7846171) for an innovative 3D printing robot design, enhancing automation in manufacturing processes.

            \end{highlights}
        \end{onecolentry}
                \vspace{0.2 cm}


\section{Patents}
\begin{onecolentry}
    % 简化日期命令，直接使用\hfill实现右对齐
    \newcommand{\addDate}[1]{\hfill\textit{#1}}
    
    \begin{highlightsforbulletentries}
    
        \item Li D, \textbf{Xue J. (2024).} \textit{Innovative 3D printing robot for automated manufacturing.} CN Patent No. 7846171. China National Intellectual Property Administration.
              
    \end{highlightsforbulletentries}
\end{onecolentry}

\section{Honors}
\begin{onecolentry}
    % 简化日期右对齐命令，直接用\hfill+斜体
    \newcommand{\addDate}[1]{\hfill\textit{#1}}
    
    \begin{highlightsforbulletentries}
        \item First-Class SMC Fellowship, \textit{Beijing SMC Education Foundation}, \textbf{Top 10\%}
              \addDate{2021} 
        \item Outstanding Student Leader Graduate of Xi’an Jiaotong University
              \addDate{2023}
        \item First Prize in the RoboMaster National Robotics Competition, \textbf{team leader in 2023}
              \addDate{2022-2024}
    \end{highlightsforbulletentries}
\end{onecolentry}

\section{Conference Presentations}
\begin{onecolentry}
    % 简化日期命令，直接使用\hfill实现右对齐
    \newcommand{\addDate}[1]{\hfill\textit{#1}}
    
    \begin{highlightsforbulletentries}
        \item Presented machine learning-based inverse design of mechanical metamaterials at the 3rd China Metamaterials Conference, \textbf{Jiaxing, China}
              \addDate{May 2024}

        \item Delivered a talk on composite material innovations at the 13th Asian-Australasian Conference on Composite Materials, \textbf{Kyoto, Japan}
              \addDate{Aug. 2024}
              
        \item Introduced recent work on inverse design and online monitoring techniques for 3D printing at the High-end Equipment Conference, \textbf{Taiyuan, China}
              \addDate{Jun. 2025}

    \end{highlightsforbulletentries}
\end{onecolentry}


\section{Internships}
\begin{onecolentry}
    % 简化日期命令，直接使用\hfill实现右对齐
    \newcommand{\addDate}[1]{\hfill\textit{#1}}
    
    \begin{highlightsforbulletentries}

        \item Structural Engineer Intern, DJI.
              \addDate{Oct. 2025 -- Jan. 2026}
               
        \item Intern, DJI (RoboMaster Department).
              \addDate{May 2024}
               
        \item Intern, Xihu Phibotnacci Company (Motor Control).
              \addDate{Jul. 2024}
    \end{highlightsforbulletentries}
\end{onecolentry}



\section{Skills}
\begin{onecolentry}
    \newcommand{\addDate}[1]{%
        \setbox0=\lastbox 
        \settowidth{\contentwidth}{\usebox0} 
        \unhbox0 
        \hfill\rlap{\textit{\makebox[\contentwidth][r]{#1}}}% 日期斜体
    }
    \begin{highlightsforbulletentries}
        \item \textbf{Language:} English (TOEFL iBT: 98), Chinese (Native)
        \item \textbf{Technical Skills:} SOLIDWORKS, ABAQUS, Python (PyTorch), STM32, MATLAB, Adobe Illustrator 
        \item \textbf{Methodologies:} Machine Learning, Finite Element Methods, 3D Printing, Robotics Control Systems
    \end{highlightsforbulletentries}
\end{onecolentry}


\end{document}
